\begin{abstract}
	Las organizaciones requieren sistemas de inteligencia artificial capaces de razonar sobre bases de conocimiento internas sin comprometer la seguridad. Sin embargo, los modelos de lenguaje grandes (LLM) no acceden de forma nativa a información confidencial, y su despliegue local suele ser inviable debido a los costos de \textit{hardware}. En este contexto, los modelos de lenguaje pequeños (SLM, 1B–13B parámetros) emergen como una alternativa viable, aunque su capacidad para soportar pipelines de generación aumentada por recuperación (RAG) con mecanismos de autorreflexión en entornos empresariales aún no está plenamente establecida. Este estudio analiza mediante un estudio de mapeo sistemático (SMS) la literatura en trabajos que emplean o integran componentes como RAG y autorreflexión aplicados al soporte de conocimiento. Siguiendo las directrices de Kitchenham y Charters (2007) y Petersen et al. (2008), se examinaron publicaciones de arXiv, NeurIPS, IEEE Xplore y ACM entre 2020 y 2025. De un total de 510 resultados iniciales, se seleccionaron 40 estudios primarios; los resultados evidencian un creciente interés en arquitecturas que combinan recuperación densa con estrategias adaptativas basadas en incertidumbre, donde un alto porcentaje de los trabajos se construyen sobre la integración de componentes existentes. Asimismo, se identifican limitaciones en la evaluación, la eficiencia y la aplicabilidad en escenarios empresariales reales. Se concluye que estas técnicas muestran potencial para mejorar el razonamiento y la adaptación de los modelos, aunque persisten desafíos para su implementación en SLM bajo restricciones de privacidad y \textit{hardware}. El estudio organiza el conocimiento existente y establece una base estructurada para futuras investigaciones.\\
	
	\noindent\textbf{Palabras clave:}\\
	RAG, Self-Reflexion, slef-induction, LLM, SLM, autonomous agent, intelligent agent
\end{abstract}

\newpage

\selectlanguage{english}

\begin{abstract}
	Organizations require artificial intelligence systems capable of reasoning over internal knowledge bases without compromising security. However, large language models (LLMs) do not natively access confidential information, and their local deployment is often unfeasible due to hardware costs. In this context, small language models (SLMs, 1B–13B parameters) emerge as a viable alternative, although their capacity to support retrieval-augmented generation (RAG) pipelines with self-reflection mechanisms in enterprise environments has not yet been fully established. This study analyzes, through a systematic mapping study (SMS), the literature on works that employ or integrate components such as RAG and self-reflection.. Following the guidelines of Kitchenham and Charters (2007) and Petersen et al. (2008), publications from arXiv, NeurIPS, IEEE Xplore, and ACM between 2020 and 2025 were examined. From an initial pool of 510 results, 40 primary studies were selected; the findings reveal a growing interest in architectures that combine dense retrieval with uncertainty-based adaptive strategies, where a high percentage of works are built upon the integration of existing components. Likewise, limitations are identified in evaluation, efficiency, and applicability in real enterprise scenarios. It is concluded that these techniques show potential for improving model reasoning and adaptation, although challenges persist for their implementation in SLMs under privacy and hardware constraints. The study organizes existing knowledge and establishes a structured foundation for future research.\\
	
	\noindent\textbf{Keywords:}\\
	RAG, Self-Reflection, self-induction, LLM, SLM, autonomous agent, intelligent agent
\end{abstract}
